
\subsubsection{2.1 Verbosity Bias and Length-Controlled Evaluation}

The problem of length bias in LLM-based evaluation is well-documented. Zheng et al. [2023] show that GPT-4 as a pairwise judge prefers longer, more structured responses even when content quality is equal. Shi et al. [2024] extend this to systematic position and length bias across judge models.

Dubois et al. [2024] address this with AlpacaEval 2.0's length-controlled win rate (LC\_winrate). Their GLM-based approach regresses response length out of pairwise comparison outcomes, yielding a debiased preference score. They report a bivariate Spearman correlation of approximately 0.94 between win\_rate and LC\_winrate, and show LC\_winrate correlates more strongly with Chatbot Arena ($\rho \approx$ 0.98). The present work extends this analysis: moving from bivariate to partial correlation, controlling for verbosity (avg\_length), reveals the capability-LC alignment strengthens to 0.985 when verbosity is properly held constant. Note that the present sample of N = 223 diverse models yields a bivariate Spearman $\rho$ = 0.966, higher than the 0.94 reported by Dubois et al., reflecting a larger and more diverse model set; the partial correlation of 0.985 exceeds this sample bivariate as well.

Hu et al. [2024] provide a mechanistic decomposition of win\_rate into a desirability component (length-independent quality) and an information mass component (length-dependent content). Their framework predicts the desirability channel survives LC correction, consistent with the finding that $\beta _win$ dominates $\beta _len$ in standardized regression.

\subsubsection{2.2 Bidirectional Human-AI Alignment}

Shen et al. [2024] provide a systematic review of bidirectional alignment between humans and AI systems, identifying a gap in empirical quantification of capability-modulated alignment asymmetry. Ji et al. [2023] survey alignment approaches broadly, noting that evaluation metric consistency across human and AI evaluators is an open problem. The present work directly addresses the empirical gap identified by Shen et al.: at population scale (N = 223), capability strongly modulates the alignment between human preference (win\_rate) and AI-debiased preference (LC\_winrate).

\subsubsection{2.3 Positioning}

This work differs from prior studies in three ways: (1) it studies the \textit{partial} relationship between capability and LC preference, explicitly controlling for verbosity; (2) it provides a Frisch-Waugh-Lovell-inspired methodological robustness check, not previously applied in this literature; and (3) it analyzes both LC\_winrate and $\Delta$ as dependent variables, providing a comparative effect-size lesson with implications for evaluation study design. These analyses complement the AlpacaEval 2.0 methodology: the results validate the LC correction as a capability-order-preserving transformation within this dataset.

